

\medskip
\noindent\SGAInlineRefTarget{sga3:iii:lemma:4-17}\textbf{Lemma 4.17.}\footnote{Editor's note: \SGARefLink{sga3:iii:lemma:4-17}{Sections 4.17} and
\SGARefLink{sga3:iii:paragraph:4-18}{4.18} have been modified to take account of the additions made in \SGARefLink{sga3:iii:paragraph:4-16}{4.16}.}
\textit{Let \(X\) be an \(S\)-group smooth over \(S\), and let
\(Y_{\mathscr J}\) be a sub-\(S_{\mathscr J}\)-group of
\(X_{\mathscr J}\), flat and locally of finite presentation over
\(S_{\mathscr J}\). For every subscheme \(Y\) of \(X\), flat over \(S\)
and lifting \(Y_{\mathscr J}\), consider the obstruction defined in
\SGARefLink{sga3:iii:definition:4-16-1}{4.16.1}:}
\[
 DY\in
 \operatorname{Hom}_{S_0}
 \bigl(Y_0\mathbin{\times_{S_0}}Y_0,N_0\bigr)
 =C^2(Y_0,N_0).
\]
\begin{enumerate}
\item[\SGAInlineRefTarget{sga3:iii:lemma:4-17:item:i}(i)] \textit{For \(Y\) to be a sub-\(S\)-group of \(X\), it is
  necessary and sufficient that \(DY=0\).}
\item[\SGAInlineRefTarget{sga3:iii:lemma:4-17:item:ii}(ii)] \textit{If \(Y_0\) is made to act on \(N_0\), by
  functoriality, through the inner automorphisms of \(Y_0\), then
  \(DY\in Z^2(Y_0,N_0)\).}
\item[\SGAInlineRefTarget{sga3:iii:lemma:4-17:item:iii}(iii)] \textit{If \(Y\) and \(Y'\) are two subschemes of \(X\),
  flat over \(S\), lifting \(Y_{\mathscr J}\) (so that the deviation
  \(d(Y,Y')\in C^1(Y_0,N_0)\) is defined; cf. \SGARefLink{sga3:iii:remark:4-5-1}{4.5.1}), then}
  \[
   DY'-DY=\partial^1d(Y,Y').
  \]
\end{enumerate}

We prove these assertions in succession.

\medskip
\noindent\SGAInlineRefTarget{sga3:iii:paragraph:4-18}\textbf{4.18. --- Proof of \SGARefLink{sga3:iii:lemma:4-17:item:i}{4.17(i)}.}
By definition, \(DY=0\) if and only if \(Y\) is stable under
the group law of \(X\). Thus \(DY=0\) if \(Y\) is a subgroup
of \(X\). Conversely, if \(DY=0\), then \(Y\) is endowed with
the induced law \(\mu_Y\), which is associative and reduces modulo
\(\mathscr J\) to the group law on \(X_{\mathscr J}\). Since \(Y\) is
flat and locally of finite presentation over \(S\), it follows from \SGARefLink{sga3:iii:proposition:3-3}{3.3}
that \(\mu_Y\) is a group law.

\medskip
\noindent\SGAInlineRefTarget{sga3:iii:paragraph:4-19}\textbf{4.19. --- Proof of \SGARefLink{sga3:iii:lemma:4-17:item:ii}{4.17(ii)}.}
We compare the two values of
\(u=c(X,Y,\mu^2\circ(i,i,i))\) computed in the following two diagrams
\((D_j)\), \(j=1,2\):

\begin{center}
 \SGAThreeNativeInput[0.94\textwidth]
   {native_diagrams/exp3/diagram_056_exp3_p063_Dj_associativity.tex}
\end{center}

Here \(f_1=(1,\pi)\), \(f_2=(\pi,1)\), and \(\mu^2\) denotes the
morphism
\[
 \mu\circ f_1=\mu\circ f_2:
 X\mathbin{\times_S}X\mathbin{\times_S}X\longrightarrow X.
\]
\footnote{Editor's note: The notation has been slightly modified, and
the beginning of the argument has been expanded.}
Put
\[
 \mu_Y=\mu\circ(i,i),\qquad
 f_j^Y=f_j\circ(i,i,i),\qquad
 \mu^{2,Y}=\mu^2\circ(i,i,i).
\]
For \(j=1,2\), let \(a_j\) and \(b_j\) be the morphisms
\[
 a_j=a_{\mu_0}\bigl((f_j^Y)_0\bigr),
 \qquad
 b_j=b_{(f_j^Y)_0},
\]
associated by \SGARefLink{sga3:iii:lemma:4-12}{Lemma 4.12} with the pair of morphisms \((f_j^Y,\mu)\).
Thus
\begin{equation}
\tag{\(\dagger\)}\label{sga3:iii:paragraph:4-19:equation:dagger}
\left\{
\begin{aligned}
 a_j:\;&
 \operatorname{Hom}_{\mathscr O_{Y_0^3}}
 \left((f_j^Y)_0^*
  \mathscr N_{Y_0\times Y_0/X_0\times X_0},
  \mathscr J\mathscr O_{Y_0^3}\right)
 \longrightarrow\\[-0.15em]
 &\qquad
 \operatorname{Hom}_{\mathscr O_{Y_0^3}}
 \left((\mu^{2,Y})_0^*\mathscr N_{Y_0/X_0},
  \mathscr J\mathscr O_{Y_0^3}\right),\\[0.35em]
 b_j:\;&
 \operatorname{Hom}_{\mathscr O_{Y_0^2}}
 \left((\mu_Y)_0^*\mathscr N_{Y_0/X_0},
  \mathscr J\mathscr O_{Y_0^2}\right)
 \longrightarrow\\[-0.15em]
 &\qquad
 \operatorname{Hom}_{\mathscr O_{Y_0^3}}
 \left((\mu^{2,Y})_0^*\mathscr N_{Y_0/X_0},
  \mathscr J\mathscr O_{Y_0^3}\right).
\end{aligned}
\right.
\end{equation}

Since
\[
 \mathscr N_{Y_0\times Y_0/X_0\times X_0}
 \simeq
 \operatorname{pr}_1^*\mathscr N_{Y_0/X_0}
 \oplus
 \operatorname{pr}_2^*\mathscr N_{Y_0/X_0}
\]
(because \(X_0\) and \(Y_0\) are flat over \(S_0\)), and
\[
 \mathscr N_{Y_0/X_0}
 \simeq
 \mathbf n_{Y_0/X_0}\otimes_{\mathscr O_{S_0}}\mathscr O_{Y_0},
\]
we have
\[
 (f_j^Y)_0^*\mathscr N_{Y_0\times Y_0/X_0\times X_0}
 \simeq
 (\mathbf n_{Y_0/X_0}\oplus \mathbf n_{Y_0/X_0})
 \otimes\mathscr O_{Y_0^3},
\]
and likewise
\[
 \begin{aligned}
 (\mu^{2,Y})_0^*\mathscr N_{Y_0/X_0}
 &\simeq \mathbf n_{Y_0/X_0}\otimes\mathscr O_{Y_0^3},\\
 (\mu_Y)_0^*\mathscr N_{Y_0/X_0}
 &\simeq \mathbf n_{Y_0/X_0}\otimes\mathscr O_{Y_0^2}.
 \end{aligned}
\]
Moreover, since \(Y_0^2\) and \(Y_0^3\) are flat over \(S_0\),
\((\dagger)\) may be rewritten as
\begin{equation}
\tag{\(\ddagger\)}\label{sga3:iii:paragraph:4-19:equation:ddagger}
\left\{
\begin{aligned}
 a_j&:\operatorname{Hom}_{S_0}(Y_0^3,N_0\oplus N_0)
       \longrightarrow\operatorname{Hom}_{S_0}(Y_0^3,N_0),\\
 b_j&:\operatorname{Hom}_{S_0}(Y_0^2,N_0)
       \longrightarrow\operatorname{Hom}_{S_0}(Y_0^3,N_0).
\end{aligned}
\right.
\end{equation}

Applying \SGARefLink{sga3:iii:proposition:4-13}{4.13} twice to \(c(X,Y,\mu^{2,Y})=u\), we obtain
\[
 \begin{aligned}
 a_1\bigl(c(X^2,Y^2,f_1)\bigr)
 +b_1\bigl(c(X,Y,\mu_Y)\bigr)
 &=u\\
 &=a_2\bigl(c(X^2,Y^2,f_2)\bigr)
 +b_2\bigl(c(X,Y,\mu_Y)\bigr).
 \end{aligned}
\]
Now \(c(X,Y,\mu_Y)=DY\), and, since
\(f_1=(1,\mu)\) and \(f_2=(\mu,1)\), we have, with the evident notation,
\[
 c(X^2,Y^2,f_1)=(0,DY),
 \qquad
 c(X^2,Y^2,f_2)=(DY,0).
\]
Consequently,
\[
 u=a_1((0,DY))+b_1(DY)
  =a_2((DY,0))+b_2(DY).
\]
The first point to note is that \(b_j\) is simply
\(\operatorname{Hom}_{S_0}((f_j^Y)_0,N_0)\), that is, the morphism
deduced from \((f_j^Y)_0\) by functoriality. The preceding identity
therefore becomes
\[
 \begin{aligned}
 a_1((0,DY))
 &-\operatorname{Hom}((\mu,1),N_0)(DY)
 +\operatorname{Hom}((1,\mu),N_0)(DY)\\
 &\hspace{9em}-a_2((DY,0))=0.
 \end{aligned}
\]
The two middle terms are the \(DY(xy,z)\) and
\(DY(x,yz)\) parts of the formula for the \(2\)-coboundary.
It remains only to identify the other two terms.

We first compute the map \(a_j\). It is obtained, by pullback along
\((f_j^Y)_0\), from the morphism of \(\mathscr O_{Y_0^2}\)-modules
\[
 P:\mathbf n_{Y_0/X_0}\otimes\mathscr O_{Y_0^2}
 \longrightarrow
 (\mathbf n_{Y_0/X_0}\oplus \mathbf n_{Y_0/X_0})
 \otimes\mathscr O_{Y_0^2}
\]
induced by multiplication in \(Y_0\). Consider the vector bundle
\(V=\mathbf V(\mathbf n_{Y_0/X_0})\). By duality, \(P\) gives
\[
 \mathbf V(P):
 V\times_{S_0}V\times_{S_0}Y_0\times_{S_0}Y_0
 \longrightarrow V\times_{S_0}Y_0\times_{S_0}Y_0,
\]
which is expressed set-theoretically by
\[
 \mathbf V(P)(u,v,a,b)=(u+\operatorname{Ad}(a)v,ab,b).
\]
This is proved exactly as is the corresponding fact for Lie algebras,
that is, for the module \(\omega^1_{Y_0/S_0}\). First observe that,
by functoriality in \(Y_0\), \(V\) is endowed with a group structure in
the category of vector bundles over \(S_0\). By the lemma already used
for Lie algebras (\SGARefLink{sga3:ii:lemma:3-10}{Expos\'e II, 3.10}), this structure coincides with the
group structure underlying its \(\mathscr O_S\)-module structure. Next,
\[
 \mathbf V\bigl(\mathbf n_{Y_0/X_0}
   \otimes_{\mathscr O_{S_0}}\mathscr O_{Y_0}\bigr)
 =\mathbf V(\mathscr N_{Y_0/X_0})
\]
is itself endowed with an \(S_0\)-group structure, namely the
semidirect product of the structure on \(V\) by that on \(Y_0\). It
remains only to identify the operations of \(Y_0\) on \(V\) to obtain
the desired formula.

We now compute the two remaining terms. Consider first
\(a_1((0,DY))\). It is computed by the following diagram, in
which \(\mathbf n\) denotes \(\mathbf n_{Y_0/X_0}\):

\begin{center}
 \SGAThreeNativeInput[0.51\textwidth]
   {native_diagrams/exp3/diagram_057_exp3_p065_module_pullback.tex}
\end{center}

Regard the vector bundles defined by these modules as schemes over
\(S_0\), and take points with values in an arbitrary scheme. If
\((u,x,y,z)\) denotes a point of
\(\mathbf V(\mathscr J)\times Y_0^3\), we obtain the following diagram:

\begin{center}
 \SGAThreeNativeInput[0.78\textwidth]
   {native_diagrams/exp3/diagram_058_exp3_p065_vector_bundle_points.tex}
\end{center}

Thus
\[
 a_1((0,DY))(x,y,z)
 =\operatorname{Ad}(x)DY(y,z),
\]
which is indeed the first term of the coboundary. Likewise,
\(a_2((DY,0))(x,y,z)=DY(x,y)\), and therefore
\[
 \begin{aligned}
 0={}&\operatorname{Ad}(x)DY(y,z)-DY(xy,z)
      +DY(x,yz)-DY(x,y)\\
   ={}&(\partial^2DY)(x,y,z).
 \end{aligned}
\]

\medskip
\noindent\SGAInlineRefTarget{sga3:iii:paragraph:4-20}\textbf{4.20. --- Proof of \SGARefLink{sga3:iii:lemma:4-17:item:iii}{4.17(iii)}.}
\footnote{Editor's note: The argument that follows has been expanded
from the original.}
We compare the two values of
\(v=c(X,Y,\mu\circ(i',i'))\) computed in the following two diagrams:

\begin{center}
 \SGAThreeNativeInput[0.78\textwidth]
   {native_diagrams/exp3/diagram_059_exp3_p066_comparison_diagrams.tex}
\end{center}

Let \(f=\mu\circ(i',i')\). Diagram \((*)\) gives
\begin{equation}
 v=DY'+f_0^*\bigl(c(X,Y,i')\bigr).
 \tag{1}\label{sga3:iii:paragraph:4-20:equation:1}
\end{equation}
Now \(Y'_0=Y_0\), and \(f_0\) is the multiplication
\(Y_0^2\to Y_0\); hence
\begin{equation}
 f_0^*\bigl(c(X,Y,i')\bigr)(x_0,y_0)
 =c(X,Y,i')(x_0y_0).
 \tag{2}\label{sga3:iii:paragraph:4-20:equation:2}
\end{equation}
Put \(c=c(X,Y,i')\). Under the identification
\(N'_0\simeq N_0\oplus N_0\), the element
\[
 c(X\times_S X,Y\times_S Y,(i',i'))
\]
identifies with the pair \((c,c)\). Thus, putting \(h=(i',i')\),
diagram \((\dagger)\) gives
\begin{equation}
 v=h_0^*(DY)+a_{\mu_0}(c,c).
 \tag{3}\label{sga3:iii:paragraph:4-20:equation:3}
\end{equation}
Since \(h_0\) is the identity map on \(Y_0^2\), we have
\(h_0^*(DY)=DY\). Finally, by the preceding
calculation of \(a_{\mu_0}\), for every \(S'\to S\) and every
\(x_0,y_0\in Y_0(S'_0)\),
\begin{equation}
 a_{\mu_0}(c,c)(x_0,y_0)
 =c(x_0)+\operatorname{Ad}(x_0)\bigl(c(y_0)\bigr).
 \tag{4}\label{sga3:iii:paragraph:4-20:equation:4}
\end{equation}
It follows that
\begin{equation}
\begin{aligned}
 (DY'-DY)(x_0,y_0)
  ={}&\operatorname{Ad}(x_0)\bigl(c(X,Y,i')(y_0)\bigr)
      -c(X,Y,i')(x_0y_0)+c(X,Y,i')(x_0)\\
  ={}&\bigl(\partial^1c(X,Y,i')\bigr)(x_0,y_0).
\end{aligned}
\tag{5}\label{sga3:iii:paragraph:4-20:equation:5}
\end{equation}
Since \(c(X,Y,i')=d(Y,Y')\) (cf. \SGARefLink{sga3:iii:proposition:4-8:item:iii}{4.8(iii)}), this proves that
\[
 DY'-DY=\partial^1d(Y,Y').
\]

\medskip
\noindent\SGAInlineRefTarget{sga3:iii:theorem:4-21}\textbf{Theorem 4.21.}
\textit{Let \(S\) be a scheme and let \(\mathscr I\) and
\(\mathscr J\) be two ideals on \(S\), such that \(\mathscr I\supset\mathscr J\) and
\(\mathscr I\mathscr J=0\). Let \(X\) be an \(S\)-group smooth over
\(S\), and let \(Y_{\mathscr J}\) be a sub-\(S_{\mathscr J}\)-group of
\(X_{\mathscr J}\), flat and locally of finite presentation over
\(S_{\mathscr J}\). Consider the \(S_0\)-functor of commutative groups
\(N_0\) defined by}
\[
 \begin{aligned}
 \operatorname{Hom}_{S_0}(T,N_0)
 ={}&\operatorname{Hom}_{\mathscr O_T}
 \left(
  \mathbf n_{Y_0/X_0}\otimes_{\mathscr O_{S_0}}\mathscr O_T,
  \mathscr J\otimes_{\mathscr O_{S_0}}\mathscr O_T
 \right),\\[-0.2em]
 &\hspace{17em}T\in\operatorname{Ob}(\mathbf{Sch}/S_0),
 \end{aligned}
\]
\textit{on which \(Y_0\) acts through the inner automorphisms of
\(X_0\).}
\begin{enumerate}
\item[(i)] \textit{For there to exist a sub-\(S\)-group of \(X\), flat
  over \(S\), which reduces to \(Y_{\mathscr J}\), it is necessary and
  sufficient that the following two conditions hold:}
  \begin{enumerate}
  \item[\((i_1)\)] \textit{There exists a subscheme \(Y\) of \(X\),
    flat over \(S\), lifting \(Y_{\mathscr J}\). This condition is
    automatically satisfied when \(Y_0\) is affine (cf. \SGARefLink{sga3:iii:proposition:4-6-5}{4.6.5}).}
  \item[\((i_2)\)] \textit{A certain canonical obstruction, an element
    of \(H^2(Y_0,N_0)\), vanishes.}
  \end{enumerate}
\item[\SGAInlineRefTarget{sga3:iii:theorem:4-21:item:ii}(ii)] \textit{If the conditions in (i) hold, the set of
  sub-\(S\)-groups \(Y\) of \(X\), flat over \(S\) and reducing to
  \(Y_{\mathscr J}\), is a principal homogeneous set under
  \(Z^1(Y_0,N_0)\).}
\end{enumerate}

Indeed, condition \((i_1)\) is necessary. Suppose that it holds, and let
\(Y\), flat over \(S\), lift \(Y_{\mathscr J}\). We must find a
\(Y'\), also flat over \(S\) and lifting \(Y_{\mathscr J}\), such that
\(DY'=0\), cf. \SGARefLink{sga3:iii:lemma:4-17:item:i}{4.17(i)}. By \SGARefLink{sga3:iii:lemma:4-17:item:iii}{4.17(iii)}, this amounts to finding
an element \(d(Y',Y)\in C^1(Y_0,N_0)\) such that
\[
 DY=\partial^1d(Y',Y).
\]
Let \(c\in H^2(Y_0,N_0)\) be the class represented by \(DY\),
which is a cocycle by \SGARefLink{sga3:iii:lemma:4-17:item:ii}{4.17(ii)}. By \SGARefLink{sga3:iii:lemma:4-17:item:iii}{4.17(iii)}, this class does not depend
on the choice of \(Y\), and its vanishing is necessary and sufficient
for the existence of \(d(Y',Y)\) satisfying the preceding equation.
This proves (i).

If \(Y\) has now been chosen so that \(DY=0\), the equation to
be solved becomes \(\partial^1d(Y',Y)=0\), which proves \SGARefLink{sga3:iii:theorem:4-21:item:ii}{(ii)}.

\medskip
\noindent\SGAInlineRefTarget{sga3:iii:remark:4-22}\textbf{Remark 4.22.}
Retain the notation of \SGARefLink{sga3:iii:theorem:4-21}{4.21}. By \SGARefLink{sga3:iii:proposition:4-15}{4.15}, \(Y_0\) is locally a complete
intersection in \(X_0\); hence \(\mathscr N_{Y_0/X_0}\) is a finite
locally free \(\mathscr O_{Y_0}\)-module, and therefore
\[
 \mathbf n_{Y_0/X_0}
 =\varepsilon_0^*(\mathscr N_{Y_0/X_0})
\]
is a finite locally free \(\mathscr O_{S_0}\)-module. Thus, writing
\[
 \mathbf n_{Y_0/X_0}^{\vee}
 =\operatorname{Hom}_{\mathscr O_{Y_0}}
  (\mathbf n_{Y_0/X_0},\mathscr O_{Y_0}),
\]
we have
\[
 \begin{aligned}
 &\operatorname{Hom}_{\mathscr O_T}
 \left(
  \mathbf n_{Y_0/X_0}\otimes_{\mathscr O_{S_0}}\mathscr O_T,
  \mathscr J\otimes_{\mathscr O_{S_0}}\mathscr O_T
 \right)\\
 &\qquad\simeq
 \mathbf n_{Y_0/X_0}^{\vee}
 \otimes_{\mathscr O_{S_0}}\mathscr J
 \otimes_{\mathscr O_{S_0}}\mathscr O_T
 \end{aligned}
\]
for every \(T\to S_0\).\footnote{Editor's note: The preceding argument
has been expanded. It shows that the following isomorphism is valid
without a flatness hypothesis. On the other hand, since \SGARefLink{sga3:iii:paragraph:4-16}{4.16} the
discussion has been restricted to \(S\)-schemes \(f:T\to S\) flat over
\(S\), in order to ensure that the group
\[
 \operatorname{Hom}_{\mathscr O_{T_0}}
 \left(
  \mathbf n_{Y_0/X_0}\otimes_{\mathscr O_{S_0}}
   \mathscr O_{T_0},
  \mathscr J\mathscr O_T
 \right),
\]
in which the obstruction \(c(X,Y,f)\) lies, coincides with
\(N_0(T_0)\); see the end of \SGARefLink{sga3:iii:paragraph:4-16}{4.16}.}
Consequently, the \(S_0\)-functor \(N_0\) is isomorphic to
\[
 \mathbf W\left(
  \mathbf n_{Y_0/X_0}^{\vee}
  \otimes_{\mathscr O_{S_0}}\mathscr J
 \right)
 \simeq
 \mathbf W\left(
  \operatorname{Hom}_{\mathscr O_{S_0}}
  (\mathbf n_{Y_0/X_0},\mathscr J)
 \right).
\]
It follows that there are isomorphisms
\[
 \begin{aligned}
 H^2(Y_0,N_0)
 &\simeq H^2\!\left(
   Y_0,\operatorname{Hom}_{\mathscr O_{S_0}}
   (\mathbf n_{Y_0/X_0},\mathscr J)\right)\\
 &\simeq H^2\!\left(
   Y_0,\mathbf n_{Y_0/X_0}^{\vee}
   \otimes_{\mathscr O_{S_0}}\mathscr J\right),\\[0.4em]
 Z^1(Y_0,N_0)
 &\simeq Z^1\!\left(
   Y_0,\operatorname{Hom}_{\mathscr O_{S_0}}
   (\mathbf n_{Y_0/X_0},\mathscr J)\right)\\
 &\simeq Z^1\!\left(
   Y_0,\mathbf n_{Y_0/X_0}^{\vee}
   \otimes_{\mathscr O_{S_0}}\mathscr J\right).
 \end{aligned}
\]

